\begin{abstract}
\sloppy
This volume unifies two sampling geometries for the Fabius--Rvachev system.
Its additive dyadic comb yields exact generating functions for each row,
Bell--Prouhet compression, terminating
Euler--Maclaurin formulae, phase-sensitive Rvachev cubature, and sharp global
interpolation instability with stable local alternatives.  On the geometric
comb $cq^{\NaturalNumbers}$ it develops Gaussian-Pascal and Jackson--Newton
calculus, exact residual and stability formulae, Mellin and regular-variation
filters, geometric-knot splines, inverse-dyadic Fabius transforms, and finite
Rvachev synthesis.  The common algebra explains why Gaussian binomials,
Bernoulli and Bell polynomials, Thue--Morse cancellation, Richardson
extrapolation, Lambert $W$, and the Rvachev sinc product repeatedly meet in the
same problem.

The text consolidates four live manuscripts and, through the additive-dyadic
manuscript, nine earlier source packages.  Repeated geometric-comb derivations
are stated once in their strongest form.  Independent extensions are retained
with their own proofs.  Every theorem, proposition, lemma, and corollary in
this volume has an explicit human-readable proof; every unproved mathematical
claim is visibly labelled as a conjecture, problem, or research question.
\end{abstract}

\begin{statusbox}{Proof and evidence terminology}
\textbf{Lean-proved} means that the theorem concordance names an exact compiled
declaration in the repository.  \textbf{Human-proved frontier result} means
that a complete manuscript proof is present here but no exact Lean declaration
has been validated.  \textbf{Conjecture} and \textbf{problem} mean that a
genuine proof obligation remains.  Numerical tables, symbolic checks, and
plots are reproducibility evidence, never proof premises.  A module-level
connection is not reported as a Lean proof unless an exact declaration has
been checked.  When compiled declarations cover only proper subclaims of a
compound manuscript result, the concordance keeps that result at manuscript
status and records the narrower Lean support without upgrading the whole row.
\end{statusbox}

\section*{Editorial scope and reading map}
\addcontentsline{toc}{section}{Editorial scope and reading map}

The geometric-comb part comes first because its affine and multiplicative
algebra supplies the reusable interpolation language.  The strongest source
is used for the finite spine: geometric Newton products, Jackson divided
differences, Gaussian-Pascal basis changes, Lagrange cardinals, complete
homogeneous residuals, and the endpoint/global stability dichotomy.  Two
independent reports then contribute only what is genuinely additional:
confluent Puiseux--log filters, Mellin symbols, regular summability, superflat
divergence, geometric-uniform laws, positive B-spline remainders, and the
reciprocal sinc-product filter.  The later parts return to the additive dyadic
comb for exact quadrature and endpoint-flat interpolation.

The provenance appendix records every source at an immutable Git revision.
The theorem concordance gives a one-to-one disposition for every source result
environment.  The source-disposition and companion-payload ledgers similarly
account for every pre-retirement file.  Git history remains the byte-for-byte
archive; the live tree contains only the canonical publication and the unique
evidence needed to reproduce it.

\subsection*{Normalizations}

Throughout, $0<q<1$ unless a result explicitly says otherwise.  A geometric
comb has nodes $cq^j$; an additive dyadic comb has mesh $2^{-m}$.  The bounded
Fabius distribution function is denoted $F$, and $\RvachevUp$ denotes
Rvachev's compactly supported even up-function.  The shared notation catalogue
imported by the master source governs $q$-Pochhammer symbols, Gaussian
coefficients, expectations, valuations, asymptotic notation, and named number
systems.  Local symbols are introduced once at first use.
